\begin{frame}{이번 주 진행과 체크리스트 — 5개 항목}
  \small
  \begin{columns}[T]
    \begin{column}{0.47\textwidth}
      \begin{tabular}{@{}p{0.8cm}p{3.9cm}p{2.4cm}@{}}
        \toprule
        순서 & 내용 & 시간 \\
        \midrule
        1 & 다른 두 팀 발표 청취(각 5~7분, 16.1 형식 — ``왜
          DL'' 질문 포함) & 50분 중 30분 \\
        2 & 체크리스트 기반 \textbf{작성 리뷰} 두 통 — 팀당 5개
          항목 + \textbf{반박 질문 1개 이상} & 발표 직후 \\
        3 & 질의·토론: 질문을 발표팀이 \textbf{그 자리}에서 답변,
          \textbf{교수가 양쪽 리뷰 품질을 함께 채점} & 나머지 시간 \\
        \bottomrule
      \end{tabular}
    \end{column}
    \begin{column}{0.51\textwidth}
      \begin{tabular}{@{}p{2.7cm}p{4.9cm}@{}}
        \toprule
        항목 & 확인할 것 \\
        \midrule
        문제 정의 & 무엇을 예측/분류하려 했는지 \textbf{명확}한가 \\
        방법론 적절성 & 데이터 특성(비정형·불균형·시퀀스)에
          \textbf{맞는 아키텍처}를 골랐는가 \\
        수식적 정당화 & 16.1 요구사항(최소 1개)을 \textbf{충족}하고
          실제로 \textbf{타당}한가 \\
        결과의 정직성 & 잘 안 된 부분을 \textbf{숨기지 않고}
          원인을 분석했는가 \\
        \textbf{``왜 DL인가''} & 16.1 필수 질문에 \textbf{구체적
          근거}로 답했는가 \\
        \bottomrule
      \end{tabular}
    \end{column}
  \end{columns}
  \vskip0.5em
  \begin{block}{8.2와 동일 + 1}
    4개 항목은 Ch8.2와 동일하되 \kb{5번째(``왜 고전 ML이 아니라
    딥러닝인가''의 타당성)}가 추가. 토론에서 발표팀이
    ``\textbf{좋은 지적입니다}''로 넘긴 질문 = 그 리뷰가
    \textbf{2점에 머물렀다}는 증거(8.2 FAQ와 동일한 논리) —
    \textbf{구체적 근거}(비정형성, Ch07 GBDT 비교)로 답했는지
    가 좋은 리뷰의 출발점.
  \end{block}
\end{frame}
